% --- scholar LaTeX fragment --------------------------------------------
% Generated by: scholar sessions export -f table
% Session: dry-principle
% \input this file from the body of a document whose
% preamble loads:
%   \usepackage{longtable}
%   \usepackage{booktabs}
%   \usepackage{array}
% ----------------------------------------------------------------------

\begingroup\footnotesize
\begin{longtable}{@{}%
  >{\raggedright\arraybackslash}p{0.44\textwidth}c%
  >{\raggedright\arraybackslash}p{0.13\textwidth}%
  >{\raggedright\arraybackslash}p{0.26\textwidth}@{}}
\caption{Träffar som rör påståendet, DRY-principen (13 citerade, 60 som stöder påståendet och 35 som kvalificerar eller motsäger det, av 715 unika träffar; de 188 angränsande och 419 felträffarna listas inte). Skälet till varje inkludering står i sista kolumnen; för maskinklassade rader anges modellens konfidens. Databaser: OpenAlex (OA; inte open access), DBLP, IEEE Xplore (IEEE), Web of Science (WoS), Scopus.}\label{tab:sok-dry}\\
\toprule
Titel & År & Leverantör & Skäl \\
\midrule
\endfirsthead
\multicolumn{4}{@{}l}{\emph{\tablename~\thetable{} (forts.)}}\\
\toprule
Titel & År & Leverantör & Skäl \\
\midrule
\endhead
\midrule \multicolumn{4}{r@{}}{\emph{forts.\ på nästa sida}}\\
\endfoot
\bottomrule
\endlastfoot
``Cloning considered harmful'' considered harmful: patterns of cloning in software & 2008 & OA & \emph{citerad}: kloning som avsiktligt verktyg \\
A Comparative Study of Software Bugs in Clone and Non-Clone Code & 2017 & OA & \emph{citerad}: klonfel vid inkonsekvent ändring \\
An Empirical Study of the Impacts of Clones in Software Maintenance & 2011 & IEEE & \emph{citerad}: forskningsplan utan resultat \\
An empirical study on inconsistent changes to code clones at the release level & 2010 & OA & \emph{citerad}: ingen effekt på utgåvenivå \\
Do Code Clones Matter? & 2009 & WoS, OA & \emph{citerad}: klonfel vid inkonsekvent ändring \\
Harmfulness of code duplication - A structured review of the evidence & 2009 & Scopus & \emph{citerad}: strukturerad översikt av beläggen \\
Harmfulness of Code Duplication - A Structured Review of the Evidence. & 2009 & DBLP & \emph{citerad}: strukturerad översikt av beläggen (dubblettpost av citerad källa) \\
Is cloned code more stable than non-cloned code? & 2008 & Scopus & \emph{citerad}: klonad kod ändras mer sällan \\
Late propagation in software clones & 2011 & IEEE & \emph{citerad}: klonfel vid inkonsekvent ändring \\
On the Relationship of Inconsistent Software Clones and Faults: An Empirical Study & 2016 & IEEE & \emph{citerad}: svag kloneffekt, utvecklarna kände till klonerna \\
Structured Review of the Evidence for Effects of Code Duplication on Software Quality & 2009 & OA & \emph{citerad}: strukturerad översikt av beläggen (utökad rapportversion av citerad källa) \\
Studying the Impact of Clones on Software Defects & 2010 & IEEE & \emph{citerad}: kloneffekten systemberoende \\
When Do You Repeat Yourself? Voices from the Trenches of Linux Kernel Maintainers on Code Duplication & 2026 & Scopus & \emph{citerad}: praktikers syn på DRY \\
A comparative study on the intensity and harmfulness of late propagation in near-miss code clones & 2016 & OA & stöder påståendet (0.98) \\
A Comparative Study on Vulnerabilities in Categories of Clones and Non-cloned Code & 2016 & IEEE & stöder påståendet (0.84) \\
A Fine-Grained Analysis on the Inconsistent Changes in Code Clones & 2020 & Scopus & stöder påståendet (0.98) \\
A Novel Approach for Predicting the Probability of Inconsistent Changes to Code Clones Based LDA & 2014 & OA & stöder påståendet (0.98) \\
An Empirical Analysis of Git Commit Logs for Potential Inconsistency in Code Clones & 2024 & WoS & stöder påståendet (0.83) \\
An Empirical Study of Code Clones from Commercial AI Code Generators & 2025 & OA & stöder påståendet (0.93) \\
An Empirical Study of Iterative Batch Refactoring for Code Smell Reduction and Software Quality Improvement & 2026 & DBLP, OA, WoS, Scopus, IEEE & stöder påståendet (0.73) \\
An empirical study on bug propagation through code cloning & 2019 & Scopus & stöder påståendet (0.99) \\
An empirical study on the fault-proneness of clone migration in clone genealogies & 2014 & IEEE & stöder påståendet (0.93) \\
An investigation of the fault-proneness of clone evolutionary patterns & 2017 & OA & stöder påståendet (0.98) \\
Analysis and removal of code clones in software product lines & 2013 & OA & stöder påståendet \\
Analyzing Clone Evolution for Identifying the Important Clones for Management & 2017 & OA & stöder påståendet (0.95) \\
Assessing Code Clone Harmfulness: Indicators, Factors, and Counter Measures & 2021 & Scopus & stöder påståendet (0.96) \\
Assessing the effect of clones on changeability & 2008 & OA & stöder påståendet (0.97) \\
Assessing the impact of bad smells using historical information & 2007 & OA & stöder påståendet \\
Associating Code Clones with Association Rules for Change Impact Analysis & 2020 & IEEE & stöder påståendet (0.91) \\
Automated Detection and Refactoring of Mock Clones in Java Projects & 2025 & WoS & stöder påståendet (0.90) \\
Bug Propagation through Code Cloning: An Empirical Study & 2017 & IEEE & stöder påståendet (0.99) \\
Cleaning up copy--paste clones with interactive merging & 2018 & Scopus & stöder påståendet (0.84) \\
Clone detection: Why, what and how? & 2010 & IEEE & stöder påståendet (0.89) \\
Code clone analysis methods for efficient software maintenance & 2006 & OA & stöder påståendet (0.91) \\
Code clone classification based on multi-dimension feature entropy & 2026 & Scopus & stöder påståendet (0.88) \\
Comparing Multiple MATLAB/Simulink Models Using Static Connectivity Matrix Analysis & 2018 & IEEE & stöder påståendet (0.87) \\
Cross-project clone consistent-defect prediction via transfer-learning method & 2023 & Scopus & stöder påståendet (0.98) \\
Detecting and characterizing semantic inconsistencies in ported code & 2013 & OA & stöder påståendet \\
Detection and analysis of near-miss software clones & 2009 & IEEE & stöder påståendet (0.74) \\
Do Code Clones Matter & 2009 & OA & stöder påståendet (0.99) \\
Efficient Code Clone Management based on Historical Analysis and Refactoring Support & 2013 & OA & stöder påståendet (0.91) \\
Eliminating Code Duplication in Cascading Style Sheets & 2017 & OA & stöder påståendet (0.88) \\
EMPIRICAL STUDIES OF CODE CLONE GENEALOGIES & 2012 & OA & stöder påståendet (avhandling bakom citerad källa) \\
Evaluating the effect of code duplications on software maintainability & 2013 & OA & stöder påståendet (0.77) \\
Evaluating the Impact of Code Smell Refactoring on the Energy Consumption of Android Applications & 2019 & OA & stöder påståendet (0.88) \\
ICCheck: A portable, language-Agnostic tool for synchronizing code clones & 2026 & Scopus & stöder påståendet (0.93) \\
Identifying clone removal opportunities based on co-evolution analysis & 2013 & Scopus & stöder påståendet (0.90) \\
Identifying Refactorable Micro-Clones through Mining Similarity Preserving Change Patterns & 2024 & IEEE & stöder påståendet (0.87) \\
Improving maintenance-consistency prediction during code clone creation & 2020 & Scopus & stöder påståendet (0.95) \\
Investigating Context Adaptation Bugs in Code Clones & 2019 & IEEE & stöder påståendet (0.99) \\
Is Late Propagation a Harmful Code Clone Evolutionary Pattern? An Empirical Study & 2021 & Scopus & stöder påståendet (0.82) \\
On the Analysis of Co-Occurrence of Anti-Patterns and Clones & 2017 & IEEE & stöder påståendet (0.95) \\
On the association between code cloning and fault-proneness: An empirical investigation & 2017 & IEEE & stöder påståendet (0.99) \\
On the characteristics of buggy code clones: A code quality perspective & 2018 & OA & stöder påståendet (0.83) \\
On the Relationships between Stability and Bug-Proneness of Code Clones: An Empirical Study & 2017 & Scopus & stöder påståendet (0.93) \\
OO in one sentence: keep it DRY, shy, and tell the other guy & 2004 & IEEE & stöder påståendet \\
Partial redesign of Java software systems based on clone analysis & 1999 & IEEE & stöder påståendet (0.93) \\
Predicting Consistency-Maintenance Requirement of Code Clonesat Copy-and-Paste Time & 2014 & IEEE & stöder påståendet (0.90) \\
Predicting Consistent Clone Change & 2016 & IEEE & stöder påståendet (0.97) \\
Ranking code clones to support maintenance activities & 2023 & Scopus & stöder påståendet (0.86) \\
Refactoring cross-project code duplication in an industrial software product line: A case study from RDW & 2025 & Scopus & stöder påståendet (0.93) \\
Relation of Code Clones and Change Couplings & 2006 & OA & stöder påståendet (bedömd på titel, sammanfattning saknas) \\
Software code cloning detection and future scope development- Latest short review & 2014 & IEEE & stöder påståendet (0.82) \\
Stop Software Rot: Keep Your Architecture DRY & 2004 & Scopus & stöder påståendet (0.77) \\
Study and Analysis of Bug Propagation in Evolving Software through Micro-clones & 2024 & IEEE & stöder påståendet (0.96) \\
The Development and Prospect of Code Clone & 2022 & OA & stöder påståendet (0.84) \\
Toxic Code Snippets on Stack Overflow & 2019 & OA & stöder påståendet (0.88) \\
Understanding Mirror Bugs in Multiple-Language Projects & 2026 & WoS & stöder påståendet (0.92) \\
Unraveling Code Clone Dynamics in Deep Learning Frameworks & 2025 & Scopus & stöder påståendet (0.88) \\
Unveiling code clones in the Eclipse IIoT software ecosystem & 2026 & Scopus & stöder påståendet (0.73) \\
Using clone detection to identify bugs in concurrent software & 2010 & IEEE & stöder påståendet (0.83) \\
Variant-Preserving Refactorings for Migrating Cloned Products to a Product Line & 2017 & WoS, Scopus & stöder påståendet (0.95) \\
Why and How to Control Cloning in Software Artifacts & 2011 & OA & stöder påståendet (0.98) \\
"Cloning considered harmful" considered harmful & 2006 & WoS, OA & kvalificerar eller motsäger påståendet (0.99) \\
A change-type based empirical study on the stability of cloned code & 2014 & Scopus & kvalificerar eller motsäger påståendet (0.96) \\
A comparative study of bug patterns in java cloned and non-cloned code & 2014 & Scopus & kvalificerar eller motsäger påståendet (0.99) \\
A study on code clone evolution analysis & 2017 & IEEE & kvalificerar eller motsäger påståendet (0.81) \\
An Empirical Study of Clone Removals & 2013 & IEEE & kvalificerar eller motsäger påståendet (0.96) \\
An Empirical Study of Fault Prediction with Code Clone Metrics & 2011 & IEEE & kvalificerar eller motsäger påståendet (0.95) \\
An empirical study of long-lived code clones & 2011 & Scopus & kvalificerar eller motsäger påståendet (0.73) \\
An Empirical Study on Inconsistent Changes to Code Clones at Release Level & 2009 & OA & kvalificerar eller motsäger påståendet (tidigare version av citerad källa) \\
An Empirical Study on the Impact of Duplicate Code & 2012 & OA & kvalificerar eller motsäger påståendet (0.98) \\
An empirical study on the maintenance of source code clones & 2009 & DBLP, OA, Scopus, IEEE & kvalificerar eller motsäger påståendet (0.98) \\
Beyond Dependencies: The Role of Copy-Based Reuse in Open Source Software Development & 2025 & WoS & kvalificerar eller motsäger påståendet (0.84) \\
Can I Clone This Piece of Code Here? & 2012 & WoS, OA & kvalificerar eller motsäger påståendet (0.98) \\
Clone stability & 2011 & Scopus & kvalificerar eller motsäger påståendet (0.99) \\
Cloned code: Stable code & 2013 & Scopus & kvalificerar eller motsäger påståendet (noterad, inte citerad) \\
Clones in deep learning code: what, where, and why? & 2022 & DBLP, OA, WoS, Scopus, IEEE & kvalificerar eller motsäger påståendet \\
Clones in Deep Learning Code: What, Where, and Why? & 2021 & OA & kvalificerar eller motsäger påståendet (0.67) \\
Code Cloning in Smart Contracts on the Ethereum Platform: An Extended Replication Study & 2023 & IEEE & kvalificerar eller motsäger påståendet (0.77) \\
Code Duplication and Reuse in Jupyter Notebooks & 2020 & OA & kvalificerar eller motsäger påståendet (0.95) \\
Code Reuse and Good Clones in Programming Education & 2025 & Scopus & kvalificerar eller motsäger påståendet (0.86) \\
Comparative stability of cloned and non-cloned code & 2012 & OA & kvalificerar eller motsäger påståendet (0.97) \\
Evolutionary Perspective of Structural Clones in Software & 2019 & IEEE & kvalificerar eller motsäger påståendet (0.80) \\
Governing the commons: code ownership and code-clones in large-scale software development & 2024 & OA & kvalificerar eller motsäger påståendet (0.76) \\
Is cloned code really stable? & 2018 & Scopus & kvalificerar eller motsäger påståendet (noterad, inte citerad) \\
Multilingual Investigation of Cross-Project Code Clones in Open-Source Software for Internet of Things Systems & 2024 & DBLP, OA, WoS, Scopus, IEEE & kvalificerar eller motsäger påståendet (0.91) \\
Refactoring patterns study in code clones during software evolution & 2017 & IEEE & kvalificerar eller motsäger påståendet (0.80) \\
Staying DRY with OO and FP & 2024 & DBLP, OA, Scopus, IEEE & kvalificerar eller motsäger påståendet (0.96) \\
Structural clones: An evolution perspective & 2018 & IEEE & kvalificerar eller motsäger påståendet (0.88) \\
Supporting the analysis of clones in software systems: A case study & 2006 & Scopus & kvalificerar eller motsäger påståendet (0.95) \\
Survey of Research on Software Clones & 2007 & OA & kvalificerar eller motsäger påståendet (0.86) \\
Toward an Understanding of Software Code Cloning as a Development Practice & 2009 & OA & kvalificerar eller motsäger påståendet (0.98) \\
Towards a Novel Conceptual Framework for Analyzing Code Clones to Assist in Software Development and Software Reuse & 2020 & Scopus & kvalificerar eller motsäger påståendet (0.86) \\
Tracking code clones in evolving software & 2008 & OA & kvalificerar eller motsäger påståendet (0.92) \\
Tracking code clones in evolving software & 2007 & WoS, OA & kvalificerar eller motsäger påståendet (0.88) \\
Visualizing and Understanding Code Duplication in Large Software Systems & 2006 & OA & kvalificerar eller motsäger påståendet (0.94) \\
Will this clone be short-lived? Towards a better understanding of the characteristics of short-lived clones & 2019 & Scopus & kvalificerar eller motsäger påståendet (0.92) \\
\end{longtable}
\endgroup
